\scanpage{105}
\[
 F_{2n}^+\xrightarrow{\sim}F_{2n+1}^+=\Gamma(\Sigma_{J},n\{t\}) \tag*{(19)}
\]
On a d'autre part $F_n^-=0$ si $n\leq 2$, et
\[
 F_3^-=\delta_2(t_E\otimes F_2)\overset{\unsure{\varepsilon}}{\simeq}t_E^{\otimes3} \tag*{(20)}
\]
supplémentaire de rang 1 de $F_2=F_3^+$ dans $F_3$.

On a alors, pour \deleted{$n\geq3$ \unsure{de|et|tout}} $n\geq1$
\[
 F_{2n+1}^-\xrightarrow{\sim}F_{2n+2}^-\simeq F_3^{-\unsure{\varepsilon}}\otimes_k F_{2n-2}^+ \tag*{(21)}
\]
Choisissons
\[
 \delta\in t_E \tag*{(22)}
\]
et un
\[
 p\in F_2\qquad \varepsilon_2(p)=\delta^{\otimes2} \tag*{(23)}
\]
posons
\[
 p'=\Theta_\delta(p)\qquad\text{donc}\quad \varepsilon_3(p')=\delta^{\otimes3} \tag*{(24)}
\]
Alors
\[
 \left\{\begin{aligned}
 &1,p\quad\text{base de}\quad F_2=F_2^+=F_3^+\\
 &p'\quad\text{base de}\quad F_3^-\\
 &1,p,p'\quad\text{base de}\quad F_3=F_3^+\oplus F_3^-
 \end{aligned}\right. \tag*{(25)}
\]
\[
 \left\{\begin{aligned}
 &\{1,p,\ldots,p^n\}\ \text{base de}\ F_{2n}^+\simeq F_{2n+1}^+\\
 &\{p',p'p\ldots,p'p^{\unsure{n-1}}\}\ \text{base de}\ F_{2n+1}^-\simeq F_{2n+2}^-
 \end{aligned}\right. \tag*{$(25^1)$}
\]
\deleted{Propos.}
\[
 F_{2n+1}^-=F_{2n+2}^-=p'F_{2n-2}^+\qquad\text{pour }\deleted{\unsure{p|p'}}\,n\in\mathbb N^*. \tag*{(26)}
\]
\uline{Proposition} 1 Considérons l'homomorphisme

$f\longmapsto f\mid J$
\[
 F_3\longrightarrow k^J
\]
Alors on a
\[
 \operatorname{Ker}\bigl(F_3\mathrel{\unsure{-|\to}} k^J\bigr)=F_3^-=k.p' \tag*{(27)}
\]
% archive sheet number: 104
